% =====================================================================
%  I. INTRODUCTION
% =====================================================================
\section{Introduction}
\label{sec:introduction}

\IEEEPARstart{A}{nyone} who uses a tiered language-model service makes an
allocation decision before typing a word. Providers sell families of models
that trade capability against price and latency, and a growing number also
expose a second control: whether, and how much, the model should reason
before it answers~\cite{openai_reasoning,anthropic_effort,google_thinking}.
Open-weight families
now ship the same two axes. Qwen3 offers dense models from 0.6B to 14B
parameters with a switch between thinking and non-thinking modes and a
controllable thinking budget~\cite{qwen3}. Reasoning traces raise accuracy
on hard problems~\cite{wei2022cot,deepseek2025r1}, but they multiply the
number of generated tokens, and the choice of when to pay for them is left
to the person holding the keyboard.

People get that choice wrong in two opposite directions. The first is
under-provisioning: start with a small, cheap configuration, notice a wrong
answer, and move up. Every failed attempt is billed in full, and so is the
prompt, which has to be resent at each step. The second is
over-provisioning. A heavyweight configuration chosen for yesterday's hard
task stays selected, and today's trivial prompts pay its premium. Reasoning
models make the second failure worse because they tend to produce long
traces for questions that need none, a pattern the efficient-reasoning
literature calls overthinking~\cite{sui2025stopoverthinking,
alomrani2025reasoningbudget}. Both failures waste tokens, time and energy,
and neither has been measured against the configuration that would have
been sufficient.

Model routing is the obvious remedy, and there is a large literature on
it~\cite{moslem2026routingsurvey}. Cascades query models in order of cost
and stop when a scorer accepts an answer~\cite{chen2023frugalgpt,
aggarwal2024automix}. Predictive routers pick a model from the query
alone~\cite{ding2024hybridllm,ong2024routellm}. More recent work routes over
reasoning mode~\cite{zhang2025synapseroute,lee2026dart} and over model and
reasoning strategy jointly~\cite{pan2026routetoreason}. These systems are
evaluated against static single-model policies, random assignment, an
oracle, or an automated cascade. None of the work we located evaluates
them against the policy people actually follow, which is manual escalation
with the prompt resent at every step. Energy-aware routers do measure power,
but they route over model choice or processing pipeline rather than
reasoning mode, and they report savings gross of the router's own
consumption~\cite{ziller2026greenserv,li2026ecothink}. Routers that control
reasoning effort report tokens and latency, not
joules~\cite{lee2026dart}.

This paper studies one-shot selection of the \emph{minimum sufficient
configuration}: the cheapest cell of a model-size by reasoning-mode grid
that still answers a prompt correctly, predicted before any generation. We
define the \emph{escalation tax} as the cost a policy incurs beyond that
configuration, and we split it into an escalation component, paid for
attempts that failed, and a stickiness component, paid for finishing on a
configuration more expensive than necessary. The study runs on hardware
that most practitioners own: a laptop with an 8\,GB consumer GPU, where
quantisation is forced, only one model fits in memory at a time, and every
joule is measured from board power rather than estimated from rated power.

The contributions are as follows.
\begin{itemize}
  \item A formal definition of the escalation tax and its decomposition,
        together with a family of manual escalation policies evaluated under
        both perfect and imperfect failure detection
        (Section~\ref{sec:methodology}).
  \item An oracle-labelled corpus in which every prompt is run through all
        eight cells of a Qwen3 size-by-mode grid, recording per-cell
        success, token counts, latency and measured energy
        (Section~\ref{sec:experimentation}).
  \item A comparison of three router families---classical classifiers,
        fine-tuned encoders, and a small language model used as a
        router---under \emph{net} accounting, in which each router's own
        tokens, latency and energy are subtracted from the saving it
        produces.
  \item A measured characterisation of the energy cost of the reasoning-mode
        axis, reported as joules per \emph{solved} query so that a cheap
        configuration that fails and forces a retry is not flattered.
  \item The same tax measured on two pools the laptop cannot host:
        thirteen public models priced by their providers, and six hosted
        cells in which reasoning effort is a paid setting, so the result is
        reported at frontier scale as well as on local hardware
        (Section~\ref{sec:results}).
  \item Two validation results that other studies on similar hardware can
        reuse: an empirical test of whether thinking and non-thinking modes
        are separate operating points at every model size, and evidence that
        a supported, monotonic NVML energy counter can nonetheless be
        non-physical on laptop GPUs.
\end{itemize}

The tax is large, and most of it comes from where a user stays rather than
how they climb. Leaving the 8B thinking cell selected pays 12.96 times the
minimum sufficient energy and answers 79.8\% of prompts; a frugal climb with
thinking off pays 0.89 times and answers 92.2\%. On these gradable tasks,
moving up a size below 8B without thinking answers more prompts than turning
thinking on, for 3.8 to 5.5 times less energy per solved prompt; on
open-ended requests, graded by a model judge that agrees with a human only
moderately, thinking is far cheaper, yet the frugal climb still pays the
least tax. How a router compares with that climb turns on how a failed route
is recovered. Sent to the largest cell, the best router spends 1.30 times
the climb's energy per solved prompt, though 64\% less than the sticky
default; recovering by climbing from the routed cell instead, every router
comes within 6\% of the climb and four are cheaper than it. Detection
matters as much: when users accept 15\% of wrong answers, the router answers
7.1 points more than the climb, because it makes fewer attempts at which to
be fooled.

The rest of the paper is organised as follows. Section~\ref{sec:literature}
surveys routing, reasoning-effort control and energy measurement for
language-model inference. Section~\ref{sec:methodology} formalises the
escalation tax and describes the routers and the cost model.
Section~\ref{sec:experimentation} gives the experimental setup, including
the corpus, labelling procedure, baselines, metrics and the two validation
gates. Sections~\ref{sec:results} and~\ref{sec:discussion} present and
interpret the results, Section~\ref{sec:future} outlines future work, and
Section~\ref{sec:conclusion} concludes.
